\documentclass[11pt,a4paper]{article}
\usepackage{amsmath,amssymb,amsthm}
\usepackage{geometry}
\usepackage{hyperref}
\usepackage{cascade-xref}        % cross-paper \xref machinery (incl. xr-hyper)
\usepackage{booktabs}
\usepackage{graphicx}
\xrhyperdoc{prelude}{cascade-series-prelude}
\xrhyperdoc{part0}{cascade-series-part0}
\xrhyperdoc{part3}{cascade-series-part3}
\geometry{margin=2.5cm}

\newtheorem{theorem}{Theorem}[section]
\newtheorem{lemma}[theorem]{Lemma}
\newtheorem{corollary}[theorem]{Corollary}
\newtheorem{definition}[theorem]{Definition}
\newtheorem{remark}[theorem]{Remark}
\newtheorem{observation}[theorem]{Observation}

\title{The Cascade Series --- Part 0a\\[0.5em]
\large How Primes Fill Spacetime: Bridge between Prelude and Part 0}
\author{RTAC}
\date{May 2026}

\begin{document}
\setlength{\emergencystretch}{3em}
\maketitle
\phantomsection
\label{paperdoc}

\begin{center}
\fbox{\parbox{0.9\textwidth}{\textbf{Status.} Structural / bridging
paper. The Prelude derives $\pi$ from primes via the Euler product
and establishes the dimension tower $\{(\mathbb{R}^d, B^d)\}_{d \in
\mathbb{N}}$; Part~0 derives the four distinguished dimensions
$\{5, 7, 19, 217\}$ and the cascade invariant $\sim 10^{-120}$ from
the Gamma-function structure on integer arguments. Between them sits
a question this paper takes up: \emph{how does the prime structure
express itself across dimensions before the distinguished-dimension
machinery of Part~0 weighs in?} The paper introduces the
prime-cascade reading of the integer cascade --- atomic ticks
(primes) and derived ticks (composites), with back-references through
the divisibility lattice --- and identifies a natural divisor-weighted
quadratic invariant $\Sigma(N) = \sum_{n \leq N} \sigma(n)
\sim (\pi^2/12) N^2$. At the current cosmic depth $N \sim 10^{60}$,
this quantity reaches $\sim 10^{120}$ --- numerically reproducing the
inverse cascade invariant of Part~0 within order-unity factors. The
paper does not derive the distinguished dimensions, the spacetime
uniqueness at $d = 4$, or the cascade's $\Lambda$ prediction
generatively from prime structure alone; those remain Parts 0 and~III
respectively. The cascade's foundational footprint stays at pure
$\mathrm{RCA}_0$ inherited from the Prelude.}}
\end{center}

\begin{abstract}
The Prelude forces $\pi$ as a convergence target from the Euler
product on primes and establishes the cascade's dimension tower over
$\mathbb{N}$. Part~0 uses Gamma-function structure on integer
arguments to derive the four distinguished dimensions $\{5, 7, 19,
217\}$ and the cascade invariant $\sim 10^{-120}$. This paper bridges
those two derivations by reading the cascade tick-by-tick in
prime-cascade terms: each Planck tick adds an integer that is either
\emph{atomic} (a prime, a fresh orthogonal axis with no relationships
to prior structure) or \emph{derived} (a composite, a lattice point
in the divisibility graph with back-references to its divisors).
Spacetime $(d \in \{1, 2, 3, 4\})$ is the cascade's
\emph{primordial era} --- the first four ticks --- with $d = 4$ as
the cascade's first derived event. After the primordial era, every
later integer carries back-references to a subset of $\{1, 2, 3, 4\}$,
and these back-references propagate recursively through the
divisibility lattice. The natural divisor-weighted content invariant
of the cascade at depth $N$ is $\Sigma(N) = \sum_{n \leq N} \sigma(n)
\sim (\pi^2/12) N^2$, with the prefactor $\pi^2/12 = \zeta(2)/2$
inheriting $\pi$ from the same Euler product the Prelude commits to.
At the current cosmic depth $N \sim 10^{60}$, $\Sigma(N) \sim
10^{120}$ --- numerically the inverse of the cascade invariant
$\Lambda \sim 10^{-120}$ that Part~0 derives via the dimensional
route through $\Omega_{19} \Omega_{217}$. Whether these two routes
are the same underlying quantity, numerically coincident but
distinct, or projections of a deeper cascade structure is the
paper's principal open question. The cascade's foundational footprint
stays at pure $\mathrm{RCA}_0$.
\end{abstract}

% ============================================================
\section{The Problem This Paper Addresses}\label{sec:problem}
% ============================================================

The Prelude derives the cascade series' starting structure --- the
dimension tower
$\{(\mathbb{R}^d, B^d)\}_{d \in \mathbb{N}}$ of finite-dimensional
Euclidean spaces with their unit balls --- from the single
pre-mathematical input $0 \neq 1$. The number-theoretic content is
forced by the Euler product on primes: $\pi$ enters the cascade as
the convergence target of $\sqrt{6\,\zeta_N(2)}$, where $\zeta_N(2)$
is the partial-computation form of $\sum 1/n^2$. The geometric
content (vector spaces, inner products, orthogonality at $\pi/2$) is
recovered from the prime-forced $\pi$ via Hurwitz's theorem and
elementary linear algebra.

Part~0 picks up from this starting structure and derives the cascade's
\emph{distinguished dimensions} $\{5, 7, 19, 217\}$ from Gamma-function
critical points on the integer axis, together with the cascade
invariant $\sim 10^{-120}$ as a product $\Omega_{19} \cdot \Omega_{217}$
of sphere areas at two of those dimensions, modified by a Gram
correction.

Between these two derivations a structural question remains open:
\emph{how does the prime structure of the Prelude express itself
across dimensions, before the Gamma-function machinery of Part~0
takes over?} The Prelude gives $\pi$; Part~0 evaluates $\Omega_d$ at
specific integer $d$. The link is mediated only through the constant
$\pi$ entering $\Omega_d = 2\pi^{d/2}/\Gamma(d/2)$. The full
arithmetic structure of $\mathbb{N}$ --- divisibility, the prime
factorisations of intermediate integers, the recursive back-reference
graph that the cascade builds tick by tick --- plays no explicit role
in this mediation.

This paper makes that arithmetic structure explicit. It reads the
cascade's tick-by-tick activity through the lens of the prime
cascade, identifies the structural events (atomic ticks, derived
ticks), exhibits the back-reference accumulation that fills the
cascade's foundational dimensions with content as cosmic time
advances, and surfaces a natural divisor-weighted quadratic invariant
that reproduces the order of magnitude of the cascade's invariant
scale at the current cosmic depth. The bridge is structural, not
generative: the distinguished dimensions and the cascade invariant
remain Part~0 content; what changes is the reading of \emph{how} the
prime structure is expressed by the time we reach Part~0.

% ============================================================
\section{From the Prelude}\label{sec:from-the-prelude}
% ============================================================

The relevant Prelude results, used here as background:

\begin{itemize}
\item Definition: $0 \neq 1$ as sole pre-mathematical input
(\xref{prelude}{def:axiom}).
\item Principle: austerity (parameter economy, minimal strength,
unobservable quotient) as the operational content of sole-input
status (\xref{prelude}{princ:austerity}).
\item Theorem: the cardinal index $d$ ranges over $\mathbb{N}$
(\xref{prelude}{thm:iterate}).
\item Theorem: the integer structure of $\mathbb{N}$ forces $\pi$ as
a convergence target via the Euler product on primes
(\xref{prelude}{thm:pi-from-primes}).
\item Theorem: the $d$ cumulated states are pairwise orthogonal in
$\mathbb{R}^d$ with $\langle e_i, e_j \rangle = \delta_{ij}$
(\xref{prelude}{thm:orth}).
\end{itemize}

\noindent The cascade's foundational footprint is pure
$\mathrm{RCA}_0$: natural numbers, computable functions on them, and
computable reals. Every analytic-looking quantity is a sequence of
rational approximations with $\mathrm{RCA}_0$-expressible convergence
behaviour. The Monotone Convergence Theorem and stronger analytic
principles serve as bookkeeping shorthand in downstream papers, not
as foundational commitments (see the Prelude status box). This paper
inherits that footprint without modification.

What the Prelude does \emph{not} establish, and what this paper
takes up:

\begin{itemize}
\item Tick-by-tick activity of the cascade beyond the existence
of the tower.
\item Structural distinction between prime ticks and composite ticks.
\item The back-reference graph that the cascade builds in the
divisibility lattice.
\item How spacetime dimensions $\{1, 2, 3, 4\}$ accumulate content
from later integer ticks.
\item A natural quadratic-in-$N$ structural invariant of the cascade
at resolution depth $N$.
\end{itemize}

% ============================================================
\section{The Prime Cascade}\label{sec:prime-cascade}
% ============================================================

Read tick by tick, the cascade adds one integer at each Planck tick:
$1, 2, 3, 4, 5, \ldots$ Each tick adds a node to the cascade's
structural record. Two kinds of node arise.

\begin{definition}[Atomic and derived ticks]\label{def:atomic-derived}
At tick $n$ the cascade adds the integer $n$ to its structure. The
tick is \emph{atomic} if $n$ is prime and \emph{derived} if $n$ is
composite (i.e., $n \geq 4$ with a non-trivial divisor). The unit
$n = 1$ is the trivial root; it is neither atomic nor derived.
\end{definition}

The structural content of a tick is determined by the divisor set
$\{d : d \mid n\}$ of the entering integer.

\begin{observation}[Atomicity is intrinsic]\label{obs:atomic-intrinsic}
Tick $n$ is atomic if and only if the only divisors of $n$ are $1$
and $n$ itself. Equivalently: atomic ticks are exactly those for
which the entering integer has no non-trivial relationship to any
prior tick. Compositeness corresponds to having at least one
non-trivial back-reference. The criterion is local to the entering
integer and requires no comparison with the rest of the cascade.
\end{observation}

\begin{definition}[Super primes, taken primes]\label{def:super-taken}
At resolution depth $N$, an atomic tick $p \leq N$ is \emph{super}
(or \emph{untaken}) if $2p > N$, and \emph{taken} if $2p \leq N$.
Equivalently: a prime $p$ is super at $N$ iff no composite multiple
of $p$ has yet entered the cascade. It is taken iff at least one
composite multiple ($2p$ being the smallest) has entered.
\end{definition}

\noindent The set of super primes at depth $N$ is the cascade's
\emph{leading edge}: orthogonal axes with no relationships in either
direction (no non-trivial divisors below; no composite multiples
above). Bertrand's postulate guarantees the leading edge is
non-empty at every $N \geq 2$: there is always at least one super
prime.

\begin{theorem}[Atomic-rate decay]\label{thm:atomic-rate}
The asymptotic density of atomic ticks at depth $N$ is $1/\ln N$
(Prime Number Theorem). The cumulative atomic fraction
$\pi(N)/N \to 0$ as $N \to \infty$.
\end{theorem}

\begin{proof}
This is the Prime Number Theorem in its standard form, provable in
$\mathrm{RCA}_0 + \mathrm{MCT}$ but with the relevant
$\mathrm{RCA}_0$-expressible content being: $\pi(N)/N$ is a
computable sequence whose convergence to zero is Cauchy in the form
$\forall\varepsilon\,\exists N$ such that $\pi(N)/N < \varepsilon$.
\end{proof}

\noindent The atomic-rate decay tells us the cascade was
atomic-dominated only briefly. The crossover from
atomic-majority to derived-majority happens at $N = 9$: at $N = 9$
the cascade has $\pi(9) = 4$ atomic ticks (primes $2, 3, 5, 7$) and
$9 - 4 - 1 = 4$ derived ticks (composites $4, 6, 8, 9$), and at
$N = 10$ derived ticks permanently outnumber atomic ones. By
$N = 100$ the cumulative atomic fraction is $25\%$; by $N = 10^6$
it is below $8\%$; at the current cosmic depth
$N \sim 10^{60}$ it is approximately $1/\ln(10^{60}) \approx 0.7\%$.

\begin{remark}[The atomic-dominated era is unbelievably brief]
The cascade spent its first nine Planck ticks in the atomic-dominated
regime --- approximately $10^{-43}$ seconds after the resolution
began. From tick~10 onward (still within $10^{-43}$ seconds) the
cascade has been derived-dominated. The remaining $\sim 10^{60}$
ticks of cosmic history correspond to derived-dominated lattice
filling at a rate that has decayed to less than one atomic event per
hundred ticks.
\end{remark}

% ============================================================
\section{The Primordial Era (Ticks 1--4)}\label{sec:primordial-era}
% ============================================================

The cascade's foundational dimensions are laid down in the first four
ticks. Each carries a specific structural role.

\begin{itemize}
\item \textbf{Tick 1: the unit.} $n = 1$ is the trivial root of the
divisibility lattice. Every later integer back-references it. The
unit is neither atomic nor derived in the sense of
Definition~\ref{def:atomic-derived}; it is the foundational element.
\item \textbf{Tick 2: first prime.} $n = 2$ is the first atomic
event. It introduces the cascade's first orthogonal axis, the only
prime that is also the additive successor of the unit.
\item \textbf{Tick 3: second prime.} $n = 3$ is the second atomic
event. Together with tick~2 it establishes the first non-trivial
arrangement of orthogonal axes.
\item \textbf{Tick 4: first derived event.} $n = 4 = 2^2$ is the
cascade's first composite. It is the first integer with a
non-trivial back-reference structure: it points to tick~2 with
weight $2$ (the only prime factor) and to tick~1 trivially. The
cascade's first lattice point.
\end{itemize}

\noindent A structural observation:

\begin{observation}[$d = 4$ as the cascade's first derived event]\label{obs:d4-first-derived}
The integer $4 = 2^2$ is the smallest composite. Its divisor lattice
is $\{1, 2, 4\}$, and it sits as the first node of the cascade's
divisibility graph above the atomic primes $\{2, 3\}$. Within the
primordial era, $d = 4$ is the unique non-atomic dimension.
\end{observation}

The observation has two immediate consequences:

\begin{itemize}
\item $d = 4$ is the cascade's first ``filling event'' --- the first
time the cascade adds a node whose role is to populate the lattice
of existing prime axes rather than to introduce a new axis.
\item $d = 4$ is the last dimension of the cascade's primordial era
before the first distinguished dimension $d = 5$ of Part~0 enters.
The dimensions $\{1, 2, 3, 4\}$ form the cascade's foundational
era, in which no specific dimension has yet been singled out by
Gamma-function critical-point machinery.
\end{itemize}

\begin{remark}[Spacetime as the primordial era]
The dimensions $\{1, 2, 3, 4\}$ are the cascade's foundational era
in the prime-cascade reading. They are also --- by the formal
derivation of \xref{part3}{paperdoc} (Lovelock uniqueness intersected
with Clifford periodicity) --- the dimensions of observable
spacetime. The two characterisations agree: spacetime is the
foundational era of the cascade. The prime-cascade observation does
not replace the Part~III derivation; it provides a structural
reading consistent with it. The observer lives at the boundary of the
primordial era, in the sense made formal by the Part~III argument.
\end{remark}

% ============================================================
\section{How Spacetime Fills with Content}\label{sec:spacetime-fills}
% ============================================================

After the primordial era ends at tick~4, every later integer
$n \geq 5$ carries back-references to a subset of $\{1, 2, 3, 4\}$
determined by its divisor structure. The back-references are
recursive: each one reaches a node that itself receives further
back-references from future ticks. The cascade does not lay down the
foundational dimensions and then leave them; it continually deposits
content into them via the back-reference graph.

\subsection*{Direct back-references}

For each $d \in \{2, 3, 4\}$, the integers $n > d$ with $d \mid n$
form an arithmetic progression of density $1/d$:
\begin{itemize}
\item Density of direct back-references to $d = 2$: $1/2$
(every even integer).
\item Density of direct back-references to $d = 3$: $1/3$
(every multiple of $3$).
\item Density of direct back-references to $d = 4$: $1/4$
(every multiple of $4$).
\end{itemize}

\noindent Per tick, the inclusion-exclusion total contribution to
spacetime $\{1, 2, 3, 4\}$ is approximately one back-reference. The
foundational dimensions are not uniformly loaded: $d = 2$ accumulates
content at twice the rate of $d = 4$.

\begin{theorem}[Cumulative direct back-references]\label{thm:direct-backref}
At depth $N$, the cumulative direct back-reference count for
spacetime dimension $d \in \{2, 3, 4\}$ is $\lfloor N/d \rfloor$,
with $\mathrm{RCA}_0$-expressible convergence under the partial-
computation reading.
\end{theorem}

\noindent The counts are: $\sim N/2$ for $d = 2$, $\sim N/3$ for
$d = 3$, $\sim N/4$ for $d = 4$. The leading factor for the total
spacetime content from direct back-references is therefore
$N \cdot (1/2 + 1/3 + 1/4) - (\text{inclusion-exclusion}) \approx
0.83\, N$.

\subsection*{Recursive back-references}

The direct back-reference count under-states the cascade's content
deposition. Each back-reference $n \to d$ passes through intermediate
divisors of $n$, and each of those intermediate nodes itself receives
future references. Counting \emph{cover-edge paths} from $n$ to $d$
in the divisibility lattice, rather than direct edges, gives the
correct content measure.

\begin{theorem}[Cover-chain count]\label{thm:chain-count}
Let $n$ be a positive integer with $d \mid n$ and write
$n/d = p_1^{a_1} p_2^{a_2} \cdots p_r^{a_r}$ in prime factorisation.
The number of maximal cover-edge chains from $n$ to $d$ in the
divisibility lattice $(\mathbb{N}, \mid)$ is the multinomial
coefficient
\[
\mathrm{chains}(n, d) \;=\; \frac{(a_1 + a_2 + \cdots + a_r)!}{a_1!\, a_2! \cdots a_r!}.
\]
\end{theorem}

\begin{proof}
A maximal chain from $n$ down to $d$ corresponds to a sequence of
prime-drops removing one factor of some $p_i$ at each step. The
total number of drops is $\Omega(n/d) = a_1 + \cdots + a_r$, and at
each step we must choose which prime to drop. The number of such
sequences, modulo equivalent reorderings of identical prime-drops,
is the multinomial coefficient stated.
\end{proof}

\noindent Worked examples:
\begin{itemize}
\item $n = 12, d = 4$: $n/d = 3$, chains $= 1!/1! = 1$. Single chain
$12 \to 4$.
\item $n = 24, d = 4$: $n/d = 6 = 2 \cdot 3$, chains $= 2!/1!1! = 2$.
Chains $24 \to 12 \to 4$ and $24 \to 8 \to 4$.
\item $n = 48, d = 4$: $n/d = 12 = 2^2 \cdot 3$, chains $= 3!/2!1!
= 3$.
\item $n = 30030 = 2 \cdot 3 \cdot 5 \cdot 7 \cdot 11 \cdot 13,\,
d = 1$: chains $= 6! = 720$. Primorials contribute factorially many
chains.
\end{itemize}

\subsection*{Asymptotic cumulative content}

\begin{theorem}[Cumulative chain content]\label{thm:chain-cumulative}
The cumulative chain-count content of spacetime dimension $d$ at
depth $N$ is
\[
T_d(N) \;=\; \sum_{\substack{m \leq N \\ d \mid m}} \mathrm{chains}(m, d) \;=\; \sum_{k \leq N/d} \mathrm{chains}(k, 1).
\]
For typical $k$ near $N$, $\mathrm{chains}(k, 1)$ has average value
asymptotically $(\ln\ln N)!\, /\, (\text{factorial corrections})$
(Erdős--Kac). The dominant contribution to $T_d(N)$ comes from
integers with many distinct prime factors --- in particular from
primorial-like integers, whose chain count grows factorially in the
number of prime factors.
\end{theorem}

\noindent The chain-count rate per tick is therefore not constant
but grows with $N$: at any cosmic epoch, the cascade is depositing
content into spacetime at a rate that exceeds the rate at any
earlier epoch. The growth is sub-exponential --- bounded by
$(\ln\ln N)!$ for typical contributions --- but unbounded as
$N \to \infty$.

\begin{remark}[Acceleration without exponentiation]
The cascade's content deposition into spacetime is not constant in
time; it accelerates. But the acceleration is slow: at $N = 10^6$ the
typical chain count per tick is approximately $7$; at $N = 10^{60}$
it is approximately $120$; at $N = 10^{1000}$ it would reach
approximately $5000$. The cascade is genuinely producing more
content per tick as cosmic time advances, but the growth is in the
$(\ln \ln N)!$ regime, far slower than de~Sitter exponential
expansion. The acceleration in content deposition is therefore a
distinct phenomenon from cosmic acceleration; this paper does not
identify them.
\end{remark}

% ============================================================
\section{The Natural Quadratic Invariant}\label{sec:quadratic-invariant}
% ============================================================

The chain count of Section~\ref{sec:spacetime-fills} counts paths
with equal weight, regardless of which divisor a back-reference
terminates at. A finer measurement uses the divisor's \emph{value}
as a weight: a back-reference to divisor $d$ contributes $d$ to the
count. This gives the cascade's \emph{divisor-weighted content}.

\begin{definition}[Divisor-weighted content]\label{def:divisor-weighted}
The divisor-weighted content of the cascade at depth $N$ is
\[
\Sigma(N) \;=\; \sum_{n \leq N} \sigma(n) \;=\; \sum_{n \leq N} \sum_{d \mid n} d.
\]
\end{definition}

\noindent Each back-reference $n \to d$ contributes weight $d$ to
$\Sigma(N)$. Small divisors are referenced often but each reference
is light-weight; large divisors are referenced rarely but each
reference is heavyweight. The two effects combine to give a
quadratic asymptotic.

\begin{theorem}[Asymptotic of divisor-weighted content]\label{thm:sigma-asymptotic}
For large $N$,
\[
\Sigma(N) \;\sim\; \frac{\pi^2}{12} N^2 \;=\; \frac{\zeta(2)}{2} N^2.
\]
\end{theorem}

\begin{proof}
Exchange the order of summation:
\[
\Sigma(N) \;=\; \sum_{d \leq N} d \cdot \lfloor N/d \rfloor.
\]
For each $d$, the contribution is $d \cdot (N/d - \{N/d\}) =
N - d\{N/d\}$. Summing over $d \leq N$ gives $N^2 - O(N \log N)$ in
naive form; the precise prefactor $\pi^2/12$ arises from the
$\zeta(2) = \pi^2/6$ identity applied via Dirichlet's method on the
average order of $\sigma(n)$, namely $\sigma(n)/n \to \pi^2/6$ on
average. Standard derivation; see any analytic number theory text.
The Cauchy form of the convergence is
$\mathrm{RCA}_0$-expressible.
\end{proof}

\noindent The prefactor $\pi^2/12 = \zeta(2)/2$ is the Euler product
on primes from the Prelude, divided by $2$. The cascade's quadratic
invariant inherits $\pi$ from exactly the same prime structure that
$\pi$ inherits from in \xref{prelude}{thm:pi-from-primes}.

\begin{observation}[Quadratic content scale]\label{obs:quadratic-scale}
At any cosmic depth $N$, the cascade's divisor-weighted content
$\Sigma(N)$ scales as $(\pi^2/12) N^2$ --- a single natural quantity
that depends on the cascade's resolution depth and the prime-forced
constant $\pi^2$, with no other inputs. At $N = 10^{60}$,
$\Sigma(N) \approx 8.2 \times 10^{119}$.
\end{observation}

% ============================================================
\section{Bridge to Part 0: Numerical Convergence at Cosmic Epoch}\label{sec:bridge-part0}
% ============================================================

Part~0 derives the cascade invariant via the dimensional route. The
relevant chain of results (\xref{part0}{thm:tower} and
\xref{part0}{sec:the-cascade-invariant}):
\begin{itemize}
\item Four distinguished dimensions $\{5, 7, 19, 217\}$ emerge as
Gamma-function critical points on the integer axis.
\item Sphere areas $\Omega_{d-1} = 2\pi^{d/2}/\Gamma(d/2)$ at these
dimensions take specific values.
\item The cascade invariant is the product
$(18/\pi^3) \cdot \Omega_{19} \cdot \Omega_{217} \cdot e^{\mathrm{Gram}}$,
which evaluates to approximately $10^{-120}$ in reduced Planck units.
\end{itemize}

\noindent The numerical match between the two routes at the current
cosmic depth:

\begin{center}
\begin{tabular}{lll}
\toprule
Route & Quantity & Value at $N \sim 10^{60}$ \\
\midrule
Prime cascade ($\Sigma^{-1}$) & $12/(\pi^2 N^2)$ & $\sim 1.2 \times 10^{-120}$ \\
Dimensional cascade & $(18/\pi^3) \Omega_{19} \Omega_{217} e^{\mathrm{Gram}}$ & $\sim 10^{-120}$ \\
\bottomrule
\end{tabular}
\end{center}

\noindent Both routes give approximately $10^{-120}$ within
order-unity factors. The prefactors $\pi^2/12$ in the prime-cascade
route and $18/\pi^3$ in the dimensional route involve different
powers of $\pi$; the match in magnitude is non-trivial.

Three readings of this convergence are available:

\begin{itemize}
\item[(a)] \emph{Same underlying quantity.} $\Sigma(N)^{-1}$ and the
dimensional cascade invariant are two expressions of the same
cascade content, prime-structural and geometric. If true, there
should be an identity at the cascade's preferred $N$ relating
$12/(\pi^2 N^2)$ to $(18/\pi^3) \Omega_{19} \Omega_{217}
e^{\mathrm{Gram}}$; the cascade's distinguished dimensions and
the cascade's resolution depth would be locked together by this
identity.
\item[(b)] \emph{Numerically coincident but distinct.} $\Sigma(N)$
depends on $N$; the dimensional cascade invariant does not. They
agree at the current cosmic epoch but diverge at other epochs. The
match would then be a property of the current $N$, possibly
anthropic.
\item[(c)] \emph{Projections of a deeper cascade structure.} Both
quantities are valid expressions of a single underlying cascade
invariant, the prime cascade giving the arithmetic form and the
dimensional cascade the geometric form, with their match at the
current cosmic epoch forced by the cascade's resolution depth being
in its mature regime.
\end{itemize}

\noindent This paper does not commit to any of the three readings.
Each is consistent with the cascade's existing predictions; each
makes different predictions about what happens at very different
$N$. Section~\ref{sec:open-questions} below states the open question
of distinguishing them.

% ============================================================
\section{The de Sitter Relation, Cascade-Form}\label{sec:de-sitter}
% ============================================================

The standard cosmological observation, restated in cascade variables:
at the current cosmic depth $N \sim 10^{60}$,
\[
\Lambda \cdot N^2 \;\sim\; 1
\]
in reduced Planck units, where $\Lambda \sim 10^{-120}$ is the
cosmological constant predicted by Part~0 and $N \sim 10^{60}$ is the
universe's age in Planck ticks (the Dirac large-number coincidence).
The Friedmann equation $\Lambda = 3 H^2$ with $H \sim 1/T \sim 1/N$
gives $\Lambda N^2 \sim 1$ at any de~Sitter-regime epoch.

The cascade's contribution to this is twofold:
\begin{itemize}
\item Part~0 supplies the specific value $\Lambda \sim 10^{-120}$
via $\Omega_{19} \Omega_{217}$.
\item This paper supplies the natural quadratic-in-$N$ structural
quantity $\Sigma(N) \sim (\pi^2/12) N^2$ whose inverse matches
$\Lambda$ at the current cosmic depth within order-unity factors.
\end{itemize}

\noindent Together these support reading~(c) of
Section~\ref{sec:bridge-part0}: the cascade's invariant scale is set
by the cascade itself, and the universe's current cosmic depth is
the value at which the cascade's two natural invariants (arithmetic
$\Sigma(N)$ and dimensional $\Omega_{19}\Omega_{217}$) numerically
agree. Whether this is an identity or a coincidence at the current
$N$ is what the open question of
Section~\ref{sec:open-questions} concerns.

% ============================================================
\section{What This Paper Does Not Do}\label{sec:does-not-do}
% ============================================================

\textbf{Does not derive the distinguished dimensions.} The set
$\{5, 7, 19, 217\}$ remains derived in \xref{part0}{thm:tower} from
Gamma-function critical points on the integer axis. This paper
observes that the first three are atomic events (primes) and the
fourth is a derived event (the composite $7 \cdot 31$), and that the
super-prime count at $N = 217$ equals $19$ (the previous distinguished
atomic event), but neither observation amounts to a prime-cascade
derivation. The Gamma-function derivation remains load-bearing.

\textbf{Does not derive spacetime stability at $d = 4$.} The formal
derivation of $d = 4$ as the observer's spacetime is in
\xref{part3}{paperdoc} (Lovelock uniqueness intersected with
Clifford periodicity). This paper observes that $d = 4$ is the
cascade's first derived event in the primordial era, and that this is
consistent with $d = 4$ being the boundary of the cascade's
foundational era before the first distinguished dimension. This is a
structural consistency observation, not a derivation.

\textbf{Does not derive $\Lambda$ from prime structure alone.} The
cascade's $\Lambda \sim 10^{-120}$ prediction continues to derive
from the dimensional route through $\Omega_{19} \Omega_{217}$ in
Part~0. The prime-cascade $\Sigma(N)$ matches this value at the
current cosmic depth, but the paper does not establish that
$\Sigma(N)^{-1}$ \emph{is} the cascade's $\Lambda$ as a function of
$N$. That identification is open.

\textbf{Does not commit to a reading.} Whether
$\Sigma(N)^{-1}$ and the dimensional cascade invariant are
the same quantity, distinct quantities that coincide at the
current cosmic epoch, or projections of a deeper structure, is
left open.

% ============================================================
\section{Open Questions}\label{sec:open-questions}
% ============================================================

\textbf{Identity question.} Is $\Sigma(N)^{-1}$ literally the
prime-cascade form of the dimensional cascade invariant? An
affirmative answer would consist of an identity at the cascade's
preferred resolution depth $N$:
\[
\frac{12}{\pi^2 N^2} \;\stackrel{?}{=}\; \frac{18}{\pi^3} \cdot \Omega_{19} \cdot \Omega_{217} \cdot e^{\mathrm{Gram}}
\]
with the cascade's preferred $N$ forced by this identity itself. The
identity, if it holds, would lock together the universe's age in
Planck ticks and the cascade's distinguished dimensions in a single
structural relation.

\textbf{Time-variation question.} Does the cascade's effective
cosmological constant vary with $N$ through $\Sigma(N)$? Reading~(a)
or~(c) of Section~\ref{sec:bridge-part0} permits or requires
variation; reading~(b) does not. The variation rate under reading~(a)
is $d\Lambda/\Lambda \sim -2 dN/N$, observationally indistinguishable
from constant $\Lambda$ at current precision but in principle
distinguishable at very high precision over cosmological timescales.

\textbf{Atomic-event explanation.} Does the cascade contain a
structural reason for the universe's age to be $\sim 10^{60}$
specifically? Under reading~(a) the cosmic depth is forced by the
identity equating $\Sigma(N)^{-1}$ and the dimensional cascade
invariant; under reading~(c) the cosmic depth is the value at
which the two cascade routes converge in their mature regime. Either
reading would constitute a cascade-internal explanation of the
Dirac large-number coincidence.

\textbf{Generative derivation of distinguished dimensions.} A
prime-structural derivation of $\{5, 7, 19, 217\}$ from austerity
alone (without going through Gamma-function critical points) would
turn the prime cascade from a structural / bridging reading into a
generative cascade. Whether this is achievable is open. The
observation that the first three distinguished dimensions are atomic
events and the fourth is a derived event with $19$ super primes at
its tick is suggestive but not yet a derivation.

% ============================================================
\section{Foundational Commitments}\label{sec:foundational-commitments}
% ============================================================

This paper inherits the Prelude's foundational footprint without
modification:

\textbf{Sole pre-mathematical input.} Definition~$0 \neq 1$
(\xref{prelude}{def:axiom}).

\textbf{Foundational footprint.} Pure $\mathrm{RCA}_0$: natural
numbers, $\Sigma^0_1$ induction, $\Delta^0_1$ comprehension,
computable functions on $\mathbb{N}$, computable reals. The Monotone
Convergence Theorem and stronger analytic principles are not
foundational commitments; they appear in this paper only as
bookkeeping shorthand when stating asymptotic results.

\textbf{Analytic content as computable processes.} The chain-count
asymptotics, $\Sigma(N) \sim (\pi^2/12) N^2$, and the divisor-density
results are all $\mathrm{RCA}_0$-expressible convergence statements
about computable sequences; no completed real is named at the limit.
At any finite resolution depth $N$ every quantity in this paper has
a specific rational value.

\textbf{Downstream bookkeeping.} The closed-form notation
(asymptotic relations, integral identities) is shorthand for
finite-$N$ statements on computable sequences, consistent with the
Prelude's framing.

\bigskip

\begin{thebibliography}{99}

\bibitem{prelude} RTAC, \extlink{\cascadebase/cascade-series-prelude.pdf}{\textit{The Cascade Series --- Prelude: Why Nothing Has Structure}} (2026).

\bibitem{part0} RTAC, \extlink{\cascadebase/cascade-series-part0.pdf}{\textit{The Cascade Series --- Part 0: The Tower}} (2026).

\bibitem{part3} RTAC, \extlink{\cascadebase/cascade-series-part3.pdf}{\textit{The Cascade Series --- Part III: General Relativity from Lovelock and Clifford}} (2026).

\end{thebibliography}

\end{document}
